\section{Teoría frente a la Realidad Empírica: Casos de Santa Cruz de Tenerife}
\label{sec:teoria_realidad}

El modelo teórico desarrollado en las secciones anteriores parte de un viario fijo del $10\%$ y de cesiones estrictamente locales ($\text{EL}+D+I = 45\,\text{m}^2\text{s}/100\,\text{m}^2\text{c}$). Para contrastar estas hipótesis con el planeamiento real de Canarias, se han auditado dos sectores de suelo urbanizable de superficie comparable ($\sim 6\text{ ha}$) del Plan General de Ordenación de Santa Cruz de Tenerife (\cite{pgou_santacruz_2005,sitcan_planeamiento_2026}): \textbf{UA-1 del SSU~5.13 Residencial El Tablero} y \textbf{SSU~5.20 Ensanche Llano del Moro~3}.

\subsection{Comparativa Cuantitativa de Desglose Superficial}

La Tabla~\ref{tab:teoria_vs_realidad} evidencia una brecha estructural cercana a \textbf{$20$ puntos porcentuales} de suelo neto residencial entre el óptimo teórico ($57{,}8\%$) y la ordenación real ejecutada o prevista ($32{,}5\%$ en El Tablero UA-1 y $25{,}9\%$ en Llano del Moro~3):

\begin{table}[htbp]
  \centering
  \caption{Contraste entre el modelo normativo teórico ($6{,}00\text{ ha}$) y dos sectores reales de Santa Cruz de Tenerife.}
  \label{tab:teoria_vs_realidad}
  \resizebox{\linewidth}{!}{%
  \begin{tabular}{@{}lccc@{}}
    \toprule
    \textbf{Magnitud del Sector} & \textbf{Modelo Teórico} & \textbf{El Tablero UA-1 (SSU 5.13)} & \textbf{Llano del Moro 3 (SSU 5.20)} \\
    \midrule
    Superficie bruta ($S_{\text{bruta}}$) & $60.000\text{ m}^2\text{s}$ ($6{,}00\text{ ha}$) & $68.337\text{ m}^2\text{s}$ ($6{,}83\text{ ha}$) & $60.263\text{ m}^2\text{s}$ ($6{,}03\text{ ha}$) \\
    Edificabilidad bruta ($\text{Gross FAR}$) & $0{,}715$ & $0{,}70$ ($E = 47.709\text{ m}^2\text{c}$) & $0{,}60$ ($E = 34.605\text{ m}^2\text{c}$) \\
    Suelo neto residencial ($\% S_{\text{bruta}}$) & $\mathbf{57{,}8\%}$ ($34.690\text{ m}^2\text{s}$) & $\mathbf{32{,}5\%}$ ($22.240\text{ m}^2\text{s}$) & $\mathbf{25{,}9\%}$ ($15.662\text{ m}^2\text{s}$) \\
    Cesiones EL / Dotacional / Parque & $32{,}2\%$ (solo SL) & $\mathbf{39{,}6\%}$ ($27.049\text{ m}^2\text{s}$) & $\mathbf{37{,}9\%}$ ($22.809\text{ m}^2\text{s}$) \\
    Viario y peatonal ($\% S_{\text{bruta}}$) & $10{,}0\%$ ($6.000\text{ m}^2\text{s}$) & $\mathbf{27{,}9\%}$ ($19.048\text{ m}^2\text{s}$) & $\mathbf{22{,}3\%}$ SGV/PS (+$13{,}9\%$ local) \\
    Reserva de Vivienda Protegida & $40\%$ de $E_{\text{teórica}}$ & $71\%$ en todo SSU~5.13 ($403$ viv.) & $35\%$ ($91$ de $260$ viv.; $10\%$ PMS) \\
    Pendiente / Cota topográfica & Terreno plano idealizado & $1{,}0\%$--$9{,}5\%$ ($368$--$406\text{ m}$) & $0{,}2\%$--$12{,}1\%$ ($524$--$568\text{ m}$) \\
    \bottomrule
  \end{tabular}%
  }
\end{table}

\subsection{Factores Físicos, Periféricos e Institucionales de Desviación}

El análisis detallado de la planimetría de ordenación y la ortofoto oficial \cite{idecanarias_wms_2026} identifica tres causas convergentes que explican por qué los estándares mínimos teóricos de viario ($10\%$) y dotaciones ($32{,}2\%$) resultan insuficientes en la práctica insular:

\begin{enumerate}
  \item \textbf{Orografía y jerarquía viaria real:} El Tablero se asienta sobre un escarpe elevado ($368$--$406\text{ m}$) con pendientes de hasta el $9{,}5\%$, y Llano del Moro alcanza el $12{,}1\%$ ($524$--$568\text{ m}$). Resolver rasantes accesibles, taludes, conectividad exterior mediante Sistemas Generales Viarios ($\text{SGV}$, $2.589\text{ m}^2\text{s}$ preexistentes en SSU~5.20) y ejes peatonales ($6.938\text{ m}^2\text{s}$ en UA-1 frente a $12.110\text{ m}^2\text{s}$ de calzada rodable) eleva el suelo viario al $22\%$--$28\%$, duplicando con creces el $10\%$ teórico.
  \item \textbf{Sobredotación por aislamiento periférico:} Al situarse alejados del centro urbano consolidado, estos sectores absorben equipamientos de barrio y sistemas generales de rango suprasectorial. En UA-1, las cesiones suman $27.049\text{ m}^2\text{s}$ frente a un mínimo legal de $19.084\text{ m}^2\text{s}$ ---un excedente de \textbf{$+7.965\text{ m}^2\text{s}$} repartido entre plazas ($14.449\text{ m}^2\text{s}$), jardines ($9.621\text{ m}^2\text{s}$), dotación docente ($1.388\text{ m}^2\text{s}$) y deportiva ($1.591\text{ m}^2\text{s}$)---. De forma paralela, Llano del Moro~3 integra $8.967\text{ m}^2\text{s}$ de \textit{S.G.~Parque Urbano} junto a $13.842\text{ m}^2\text{s}$ de sistemas locales, y aloja explícitamente el $10\%$ de PMS en su parcela \textit{Residencial VP}. Dentro de UA-1 ($20.899\text{ m}^2\text{s}$ VPO y $1.341\text{ m}^2\text{s}$ libre \textit{C-5}), las manzanas \textit{C-2} ($3.110$ y $3.111\text{ m}^2\text{s}$) corroboran además cómo las ordenanzas reales restringen la ocupación y altura de la vivienda libre frente al mayor volumen protegido ($80.778\text{ m}^2\text{c}$ y $112.015$--$112.021\text{ m}^2\text{s}$ en todo el SSU~5.13).
  \item \textbf{Incertidumbre jurídica y parálisis de ejecución:} La auditoría previa del planeamiento de Santa Cruz identificó el SSU~5.13 (\textit{Actuación~1}) como caso testigo de bloqueo inversor: pese a segregarse UA-1 (ejecución privada) de las fragmentadas UA-2 y UA-3 (cooperación) para acelerar su gestión, la anulación judicial del PGO de 2013 mediante la Sentencia del TSJC~116/2014 \cite{tsjc_116_2014}, la posterior nulidad de la Ordenanza Provisional de 2022 (Sentencia TSJC~000070/2023 \cite{tsjc_070_2023}) y el vencimiento de los plazos del horizonte $2015$--$2023$ provocaron que las manzanas urbanizadas permanezcan vacías sobre tierra más de una década después (Figuras~\ref{fig:fig8_real_sc} y \ref{fig:fig_ge_3d}).
\end{enumerate}

\begin{figure}[htbp]
  \centering
  \includegraphics[width=0.88\linewidth]{img/fig8_real_sc_tenerife_case.pdf}
  \caption{Contrastación directa entre el reparto superficial del modelo teórico óptimo y la ordenación pormenorizada sobre ortofoto oficial en El Tablero (UA-1, SSU~5.13) y Llano del Moro~3 (SSU~5.20).}
  \label{fig:fig8_real_sc}
\end{figure}

\begin{figure}[htbp]
  \centering
  \includegraphics[width=0.95\linewidth]{img/el_tablero_google_earth_3d.png}
  \caption{Validación tridimensional sobre el escarpe de El Tablero (UA-1): coincidencia exacta entre los polígonos de zonificación del PGO, la red viaria ejecutada y los solares residenciales sin edificar tras la parálisis urbanística.}
  \label{fig:fig_ge_3d}
\end{figure}
